\subsection{Budget-neutral swaps under model uncertainty}
\label{sec:swaps}

The policy trees in Section~\ref{sec:results} may change substantially when the
resilience scores change. A funder that already has a vulnerability ranking can make a
more limited decision. It can ask whether available estimates agree that replacing
some countries on that list with others would improve welfare. The criterion follows
the idea of measuring robust improvement relative to a reference policy
\citep{kalluszhou2021}. Here the candidate class exchanges a fixed number of places,
and the uncertainty set comes from alternative estimated heat-sensitivity scores.
It does not encode a sensitivity model for unobserved confounding.

\subsubsection{Allocation rule}

Let $q_j$ be the number of places available in a batch of $n_j$ eligible countries
and let $\pi_0$ select the $q_j$ countries with the largest equity-weighted
exposure, $w_iH_i(1)$ from Equation~\eqref{eq:coverage}. The reference makes no use of
yield outcomes. We build a library of swap rules from training countries. Each
rule specifies a feature, a cutoff, the preferred side of the cutoff and a cap
on replacements. Cutoffs are feature deciles in the training sample. The cap is
one quarter, one half or all of the available places, rounded down.

At assignment time a rule admits the highest-priority unselected countries on
its preferred side. It drops an equal number of the lowest-priority selected
countries on the other side. The number of replacements is bounded by the cap
and by the number of countries available on either side of the exchange. This
construction preserves the quota in every batch. It can use the same fitted
rule when the mix of eligible countries changes, provided their features and
reference priorities are available. Candidate cutoffs and replacement fractions
are selected on training data, and quota filling uses only features in the
assignment batch. We include the reference itself as the no-swap candidate.

Write $B_{is}$ for country $i$'s equity-weighted resilience score under model
specification $s$. For the training batch, define
\begin{equation}\label{eq:swap-criterion}
 G_s(\pi) = \frac{1}{n_j}\sum_{i=1}^{n_j}
             [\pi_i-\pi_{0i}] B_{is},
 \qquad
 \hat\pi = \arg\max_{\pi\ \text{in library}} \min_s G_s(\pi).
\end{equation}
We evaluate all candidates and keep the reference when the best worst-case gain
is zero. The comparison is paired within model: each candidate and the reference
are valued using the same scores. Subtracting the reference separately in each
model avoids letting one model's low absolute welfare dictate the choice when
that model still favours a change to the allocation.

\paragraph{Financing constraint and guarantee.}
The two lists in each exchange have the same length, so
$\sum_i\pi_i=\sum_i\pi_{0i}=q_j$. Since the reference is feasible and has
$G_s(\pi_0)=0$ for every $s$, the chosen rule has $\min_s G_s(\hat\pi)\geq 0$
on the training batch. If the score vector were any convex combination of the
scenario vectors, its paired gain would be the same convex combination of
$G_s(\hat\pi)$ and hence non-negative there. This is a guarantee about a specified
score set and batch. Sampling error, common model bias and a different
assignment batch can all reverse the comparison in realised welfare.

\subsubsection{Evaluation design}

We form three scenario vectors from the same outcome data
($\epsilon=1$ throughout). Two use orthogonal augmented scores with propensity
clipping at 0.10 or 0.05 and the third uses the plug-in slope difference. The
benchmarks maximise the training gain from the 0.10-clipped DR score, the plug-in
score, or the mean of the three scenario scores over the identical library. This makes the number of places,
feature cutoffs, swap caps and reference common across methods. We also report
the unchanged exposure reference.

\paragraph{Simulation.}
Each replication independently draws 150 training and 300 test countries, each
with six observations. Four standard normal country covariates represent income,
connectivity, agricultural dependence and climate. Exposure increases with
agricultural dependence and falls with income. Treatment propensity is
$\operatorname{logit}^{-1}\{a(\mathrm{income}+0.7\,\mathrm{connectivity})\}$,
with $a=1$ for moderate overlap and $a=3$ for weak overlap. The heat shock is
heteroskedastic. In the first two designs the true moderation effect increases
at a connectivity threshold and falls in countries that are both warmer and
more agricultural. The third design uses moderate overlap and a common effect
of 0.30, so the reference is optimal. Training scores are formed by three-fold
country cross-fitting. We evaluate the selected rule on the independent test
countries using known welfare. Regret is the difference from the best member
of the training-derived library on that test draw, including the reference.
All comparisons use the same draws within each design.

\paragraph{Agriculture.}
We apply the algorithm to the 126-country policy sample. We split countries
into five folds. For each fold, all nuisance fitting, formation of training
scores and rule selection exclude its countries. Training scores use inner
country folds. The shock and outcome residuals used to fit preliminary slopes
are formed by a further split within the available fitting countries.
We score the held-out countries using a DR score whose nuisance functions are
trained outside their fold. We use clipping at 0.10 for evaluation and repeat
evaluation at 0.05. The rules see only features and priority when assigning
held-out places. We allocate quotas proportionally across batches with
cumulative rounding so the total is 37 for every method. As a geographic
stress test, we repeat the exercise leaving out entire World Bank regions.
Percentile intervals resample country contributions with the learned rules
and scores held fixed. They describe sampling variation conditional on those
fits and do not account for spatial dependence or model uncertainty.
